\section{Discussion and Conclusion}

We now discuss the implications of~\Md~on the underlying learning dynamics of LLMs. Long-term poisoning attacks on language models are not new. For example, we saw the creation of \textit{click, content, and troll} farms -- a form of human `language models', whose job is to misguide social networks and  search algorithms. %
The negative effect  these poisoning attacks had on search results led to changes in search algorithms: \eg, Google downgraded farmed articles\footnote{\url{https://googleblog.blogspot.com/2011/02/finding-more-high-quality-sites-in.html}}, putting more emphasis on content produced by trustworthy sources \eg~education domains, while DuckDuckGo removed them altogether\footnote{\url{https://www.technologyreview.com/2010/07/26/26327/the-search-engine-backlash-against-content-mills/}}.  

What is different with the arrival of LLMs is the scale at which such poisoning can happen once it is automated. Preserving the ability of LLMs to model low-probability events is essential to the fairness of their predictions: such events are often relevant to marginalised groups. %
Low-probability events are also vital to understand complex systems~\citep{taleb2007black}. %

Our evaluation suggests a ``first mover advantage'' when it comes to training models such as LLMs. In our work we demonstrate that training on samples from another generative model can induce a distribution shift, which over time causes~\Md. This in turn causes the model to mis-perceive the underlying learning task. To make sure that learning is sustained over a long time period, one needs to make sure that access to the original data source is preserved and that additional data not generated by LLMs remain available over time. %
The need to distinguish data generated by LLMs from other data raises questions around the provenance of content that is crawled from the Internet: it is unclear how content  generated by LLMs can be tracked at scale. 
One option is community-wide coordination to ensure that different parties involved in LLM creation and deployment share the information needed to resolve questions of provenance. 
Otherwise, it may become increasingly difficult to train newer versions of LLMs without access to data that was crawled from the Internet prior to the mass adoption of the technology, or direct access to data generated by humans at scale.








